
\documentclass[11pt,letterpaper,oneside]{report}
\usepackage[T1]{fontenc}
\usepackage[utf8]{inputenc}
\usepackage{lmodern}
\usepackage{inconsolata}
\usepackage{geometry}
\geometry{margin=0.89in,headheight=16pt,headsep=19pt,footskip=26pt}
\usepackage{xcolor}
\definecolor{navy}{HTML}{173758}
\definecolor{teal}{HTML}{157B89}
\definecolor{pale}{HTML}{E9F0F3}
\definecolor{codegray}{HTML}{F5F7F9}
\definecolor{muted}{HTML}{667785}
\usepackage{booktabs,longtable,tabularx,array,calc}
\usepackage{amsmath,amssymb}
\usepackage{graphicx}
\usepackage{microtype}
\usepackage{enumitem}
\usepackage{fancyhdr}
\usepackage{titlesec}
\usepackage{listings}
\usepackage{upquote}
\usepackage{xurl}
\usepackage{hyperref}
\hypersetup{colorlinks=true,linkcolor=navy,urlcolor=teal,bookmarksopen=true,pdfauthor={Software Build Engineering Documentation Suite}}
\urlstyle{same}
\setcounter{tocdepth}{2}
\setcounter{secnumdepth}{2}
\setlength{\parindent}{0pt}
\setlength{\parskip}{5pt plus 1pt}
\setlength{\emergencystretch}{3em}
\setlist[itemize]{leftmargin=*,topsep=4pt,itemsep=2pt}
\setlist[enumerate]{leftmargin=*,topsep=4pt,itemsep=2pt}
\titleformat{\chapter}[display]{\normalfont\bfseries\color{navy}}{\fontsize{11}{12}\selectfont CHAPTER \thechapter}{7pt}{\fontsize{23}{28}\selectfont}
\titlespacing*{\chapter}{0pt}{-6pt}{19pt}
\titleformat{\section}{\Large\bfseries\color{navy}}{\thesection}{0.65em}{}
\titleformat{\subsection}{\large\bfseries\color{teal}}{\thesubsection}{0.65em}{}
\pagestyle{fancy}
\fancyhf{}
\fancyhead[L]{\small\textcolor{muted}{PROFESSIONAL CMAKE ENGINEERING}}
\fancyhead[R]{\small\textcolor{muted}{\nouppercase{\leftmark}}}
\fancyfoot[C]{\small\textcolor{muted}{\thepage}}
\renewcommand{\headrulewidth}{0.3pt}
\fancypagestyle{plain}{\fancyhf{}\fancyfoot[C]{\small\thepage}\renewcommand{\headrulewidth}{0pt}}
\lstdefinelanguage{cmake}{sensitive=false,morekeywords={cmake_minimum_required,project,add_library,add_executable,target_compile_features,target_include_directories,target_link_libraries,target_sources,target_compile_options,target_compile_definitions,target_link_options,add_subdirectory,include,if,else,elseif,endif,foreach,endforeach,function,endfunction,macro,endmacro,set,option,find_package,find_program,add_custom_command,add_custom_target,install,configure_file,configure_package_config_file,write_basic_package_version_file,set_target_properties,set_tests_properties,add_test,message,FetchContent_Declare,FetchContent_MakeAvailable,return,cmake_parse_arguments},morecomment=[l]{\#},morestring=[b]"}
\lstdefinelanguage{sh}{morekeywords={cmake,ctest,cpack,find,sha256sum,clang-tidy,echo,cp,make,ninja,git},morecomment=[l]{\#},morestring=[b]"}
\lstdefinelanguage{yaml}{morecomment=[l]{\#},morestring=[b]",morekeywords={name,on,permissions,jobs,steps,uses,run,with,env,matrix,include,if}}
\lstdefinelanguage{text}{}
\lstset{basicstyle=\ttfamily\footnotesize,keywordstyle=\color{navy}\bfseries,commentstyle=\color{muted},stringstyle=\color{teal},backgroundcolor=\color{codegray},frame=single,rulecolor=\color{pale},breaklines=true,breakatwhitespace=false,showstringspaces=false,columns=fullflexible,keepspaces=true,tabsize=2,xleftmargin=5pt,xrightmargin=5pt,aboveskip=9pt,belowskip=9pt}
\providecommand{\tightlist}{\setlength{\itemsep}{0pt}\setlength{\parskip}{0pt}}
\providecommand{\passthrough}[1]{#1}
\providecommand{\real}[1]{#1}

\begin{document}
\hypersetup{pageanchor=false}\begin{titlepage}
\color{navy}
\vspace*{1.15in}
{\Large\bfseries PROFESSIONAL CMAKE\\[7pt]SOFTWARE BUILD ENGINEERING\par}
\vspace{0.32in}\noindent\textcolor{teal}{\rule{0.9\linewidth}{2.3pt}}\par
\vspace{0.45in}
{\parbox{0.95\textwidth}{\raggedright\fontsize{24}{29}\selectfont\bfseries Build Architecture and Design Patterns\par}}
\vspace{0.3in}
{\Large\color{teal}Views, patterns, decisions, trust boundaries\par}
\vfill
{\large DOCUMENT 04 OF 08\par}
\vspace{0.1in}
{\small Companion synthesis to Craig Scott, \emph{Professional CMake: A Practical Guide}, 20th Edition (2025).\\
Original learning/engineering material. Not affiliated with or endorsed by the author.\par}
\vspace{0.2in}
{\small 17 October 2026\par}
\end{titlepage}
\hypersetup{pageanchor=true}\clearpage\pagenumbering{roman}\tableofcontents\clearpage\pagenumbering{arabic}

\textbf{Mandate.} Model CMake as an executable product architecture
rather than a collection of scripts. \textbf{{[}BOOK-GROUNDED{]}}
Targets, usage requirements, superbuilds, exported configs, presets and
project structure map to Chapters 4, 16, 18--19, 34, 38--43.
\textbf{{[}ENGINEERING EXTENSION{]}} Architecture viewpoints, design
decision records, supply-chain boundaries and evaluation criteria are
proposed enterprise practices.

\chapter{Architecture viewpoints and
stakeholders}\label{architecture-viewpoints-and-stakeholders}

\begin{longtable}[]{@{}
  >{\raggedright\arraybackslash}p{(\columnwidth - 6\tabcolsep) * \real{0.2500}}
  >{\raggedright\arraybackslash}p{(\columnwidth - 6\tabcolsep) * \real{0.2500}}
  >{\raggedright\arraybackslash}p{(\columnwidth - 6\tabcolsep) * \real{0.2500}}
  >{\raggedright\arraybackslash}p{(\columnwidth - 6\tabcolsep) * \real{0.2500}}@{}}
\toprule\noalign{}
\begin{minipage}[b]{\linewidth}\raggedright
Viewpoint
\end{minipage} & \begin{minipage}[b]{\linewidth}\raggedright
Stakeholder
\end{minipage} & \begin{minipage}[b]{\linewidth}\raggedright
Concern
\end{minipage} & \begin{minipage}[b]{\linewidth}\raggedright
Required architectural evidence
\end{minipage} \\
\midrule\noalign{}
\endhead
\bottomrule\noalign{}
\endlastfoot
Source organization & component developer & which
\passthrough{\lstinline!CMakeLists.txt!} owns a subsystem? & source tree
and declared target inventory \\
Target graph & build engineer & compile/link requirement propagation &
directed dependency graph \\
Configuration & release engineer & presets and toolchain compatibility &
environment-to-artifact matrix \\
Distribution & SDK consumer & installed exported target behavior &
fresh-consumer test and ABI policy \\
Security / provenance & AppSec / supply-chain engineer & who can execute
code during configure/build? & trust-zone model, dependency inventory,
attestations \\
Runtime & operator & loader, path, plugins, resource dependencies &
installed runtime verification \\
\end{longtable}

\section{Target graph and edge
semantics}\label{target-graph-and-edge-semantics}

\begin{lstlisting}
             BuildSuite::core  (PUBLIC C++17, headers)
                    ^                     ^
                    | PRIVATE link        | PRIVATE link
                 bsuite_cli             bsuite_tests
                    |                     |
               runtime CLI            CTest case

  Install context: bsuite_core -> BuildSuiteTargets.cmake
                              -> BuildSuiteConfig.cmake
                              -> external find_package consumer
\end{lstlisting}

\textbf{Interpretation.} A \passthrough{\lstinline!PRIVATE!} link from
CLI to the library does not imply library requirements become private to
the compiler: the library's \passthrough{\lstinline!INTERFACE\_*!}
settings still supply the consuming CLI's compile/link context. The
exported target is an installation contract; it is not another compiled
binary.

\chapter{Pattern catalog}\label{pattern-catalog}

\section{Pattern P1: Target-owned
requirements}\label{pattern-p1-target-owned-requirements}

\textbf{Problem:} Global include directories and flags pollute unrelated
targets, particularly when a repository is consumed with
\passthrough{\lstinline!add\_subdirectory!}. \textbf{Solution:} assign
include directories, definitions, features and link requirements to each
target; use \passthrough{\lstinline!PUBLIC!},
\passthrough{\lstinline!PRIVATE!}, and
\passthrough{\lstinline!INTERFACE!} based on the consumer contract.
\textbf{Verification:} link an independent target without extra include
or linker flags.

\begin{lstlisting}
target_include_directories(core PUBLIC
  $<BUILD_INTERFACE:${CMAKE_CURRENT_SOURCE_DIR}/include>
)
target_compile_features(core PUBLIC cxx_std_17)
\end{lstlisting}

\textbf{Anti-pattern:}
\passthrough{\lstinline!include\_directories(...)!} at repository root
used as a substitute for declaring the library's API.

\section{Pattern P2: Interface contract
aggregation}\label{pattern-p2-interface-contract-aggregation}

\textbf{Problem:} Multiple first-party targets need the same optional
warnings or instrumentation. \textbf{Solution:}
\passthrough{\lstinline!INTERFACE!} libraries as opt-in policies. A
target must explicitly link a policy; do not export developer
diagnostics as consumer requirements.

\begin{lstlisting}
add_library(project_warnings INTERFACE)
target_compile_options(project_warnings INTERFACE
  $<$<CXX_COMPILER_ID:GNU,Clang>:-Wall;-Wextra>
)
# Only for first-party build targets:
target_link_libraries(bsuite_core PRIVATE project_warnings)
\end{lstlisting}

\textbf{Caveat:} Check option/list quoting and platform support;
diagnostics \passthrough{\lstinline!-Werror!} should not be a public ABI
contract.

\section{Pattern P3: Thin root, owned
subdirectories}\label{pattern-p3-thin-root-owned-subdirectories}

\textbf{Problem:} Top-level files become non-local, order-sensitive
state machines. \textbf{Solution:} top-level declares project, policy
and feature gates; subdirectories own related targets. Expose reusable
functionality via targets and functions, not implicit mutable global
variables.

\begin{lstlisting}
CMakeLists.txt           project and options
cmake/                   reusable helper modules
components/network/      network targets and local tests
components/storage/      storage targets and local tests
apps/                    executable composition
tests/integration/        cross-component tests
\end{lstlisting}

\textbf{Review:} Can one component be used as a child project without
overriding a parent toolchain, build type, linker or install prefix?

\section{Pattern P4: Installable package
boundary}\label{pattern-p4-installable-package-boundary}

\textbf{Problem:} A library works in a monorepo but cannot be consumed
as an SDK. \textbf{Solution:} install headers and binaries; export
namespace targets; create a relocatable config/version file; test from
an independent build. \textbf{Evidence:}
\passthrough{\lstinline!find\_package!} succeeds after relocation;
installed config does not contain source/build paths.

\section{Pattern P5: Controlled dependency
ingress}\label{pattern-p5-controlled-dependency-ingress}

\textbf{Problem:} Network-fetched dependency revisions drift, and
configure-time CMake scripts may execute arbitrary commands.
\textbf{Solution:} define one acquisition policy per environment:
validated package config, organization-controlled provider, pinned
FetchContent source or isolated ExternalProject. \textbf{Review:} does
the release build work without fetching mutable content? What integrity
data and license metadata are available?

\section{Pattern P6: Presets as supported
products}\label{pattern-p6-presets-as-supported-products}

\textbf{Problem:} Build knowledge exists only in developers' terminal
history. \textbf{Solution:} commit a small, meaningful
\passthrough{\lstinline!CMakePresets.json!}: Debug, Release, validation;
let developers customize
\passthrough{\lstinline!CMakeUserPresets.json!}. Each preset is a
declared combination of generator, architecture, binary directory, cache
and toolchain assumptions. \textbf{Test:} enumerate and execute
supported presets from a fresh checkout.

\section{Pattern P7: Hermetic-ish verification
environment}\label{pattern-p7-hermetic-ish-verification-environment}

\textbf{Problem:} Passing tests depend on undeclared compilers, paths,
package registries or host resources. \textbf{Solution:} explicitly
record toolchain versions and dependencies, run offline smoke jobs,
control environment variables and use fresh build directories.
\textbf{Caveat:} hermeticity is not guaranteed by CMake; the build
environment and dependencies must enforce it.

\section{Pattern P8:
Instrumentation-as-variant}\label{pattern-p8-instrumentation-as-variant}

\textbf{Problem:} Sanitizer and coverage flags contaminate release
targets and cached build trees. \textbf{Solution:} dedicated
presets/build directories with compatible compile/link instrumentation
and controlled test subsets. \textbf{Test:} produce a known sanitizer
failure; demonstrate Release remains uninstrumented.

\section{Pattern P9: Native packaging
adapter}\label{pattern-p9-native-packaging-adapter}

\textbf{Problem:} One artifact layout does not meet all OS distribution
requirements. \textbf{Solution:} first define the install contract; add
generator-specific CPack configurations and external
signing/notarization steps. \textbf{Review:} sign the final artifact,
then verify signature, runtime paths and package content on the
receiving platform.

\chapter{Pattern selection matrix}\label{pattern-selection-matrix}

\begin{longtable}[]{@{}
  >{\raggedright\arraybackslash}p{(\columnwidth - 6\tabcolsep) * \real{0.2500}}
  >{\raggedright\arraybackslash}p{(\columnwidth - 6\tabcolsep) * \real{0.2500}}
  >{\raggedright\arraybackslash}p{(\columnwidth - 6\tabcolsep) * \real{0.2500}}
  >{\raggedright\arraybackslash}p{(\columnwidth - 6\tabcolsep) * \real{0.2500}}@{}}
\toprule\noalign{}
\begin{minipage}[b]{\linewidth}\raggedright
Criterion
\end{minipage} & \begin{minipage}[b]{\linewidth}\raggedright
One normal build
\end{minipage} & \begin{minipage}[b]{\linewidth}\raggedright
Superbuild with ExternalProject
\end{minipage} & \begin{minipage}[b]{\linewidth}\raggedright
Organization dependency provider
\end{minipage} \\
\midrule\noalign{}
\endhead
\bottomrule\noalign{}
\endlastfoot
Shared target graph & Strong & Weak across external projects & Depends
on provider \\
Toolchain isolation & Limited & Strong & Provider-dependent \\
Dependency build lifecycle & Single configure/build & Independent
configure/build/install & Acquisition delegated \\
Source-level code generation & Straightforward & Requires handoffs &
Depends on acquired package \\
Best fit & cohesive components & SDK orchestration / incompatible
dependencies & centrally governed package sourcing \\
Governance concern & untrusted source CMake & prebuilt artifact
provenance & injection file change control \\
\end{longtable}

Do not make superbuild selection merely on project size. Choose
according to toolchain isolation, binary compatibility, dependency
lifetimes, organization of install interfaces, and CI caching costs.

\chapter{Quality attribute scenarios}\label{quality-attribute-scenarios}

\begin{longtable}[]{@{}
  >{\raggedright\arraybackslash}p{(\columnwidth - 6\tabcolsep) * \real{0.2500}}
  >{\raggedright\arraybackslash}p{(\columnwidth - 6\tabcolsep) * \real{0.2500}}
  >{\raggedright\arraybackslash}p{(\columnwidth - 6\tabcolsep) * \real{0.2500}}
  >{\raggedright\arraybackslash}p{(\columnwidth - 6\tabcolsep) * \real{0.2500}}@{}}
\toprule\noalign{}
\begin{minipage}[b]{\linewidth}\raggedright
Quality
\end{minipage} & \begin{minipage}[b]{\linewidth}\raggedright
Stimulus
\end{minipage} & \begin{minipage}[b]{\linewidth}\raggedright
Measurement
\end{minipage} & \begin{minipage}[b]{\linewidth}\raggedright
Architectural response
\end{minipage} \\
\midrule\noalign{}
\endhead
\bottomrule\noalign{}
\endlastfoot
Portability & add Windows runner & configure + build + install + consume
succeeds & isolate OS-specific configuration \\
Maintainability & replace implementation module & unrelated targets
unaffected & local target ownership \\
Build performance & edit one public header & timed rebuild and target
fan-out & minimize public header exposure \\
Reproducibility & clean build on isolated runner & compare provenance
and output hashes where deterministic & version and control
toolchains/deps \\
Security & dependency registry compromised & build blocked or provenance
discrepancy reported & pinned verified packages / policy gate \\
Developer experience & new machine onboarding & time to first passing
CTest & small portable presets and clear diagnostics \\
\end{longtable}

\chapter{Architecture decisions to
document}\label{architecture-decisions-to-document}

\section{ADR-001: Target graph versus global
flags}\label{adr-001-target-graph-versus-global-flags}

\textbf{Context:} multiple libraries, executable entry points,
Windows/Linux consumers. \textbf{Decision:} target-local usage
requirements. \textbf{Alternatives:} directory-global variables and
flags. \textbf{Consequences:} more explicit declarations; increased
reuse, transparency and auditability. \textbf{Verification:} independent
component and install-consumer tests.

\section{ADR-002: Dependency
acquisition}\label{adr-002-dependency-acquisition}

\textbf{Context:} CI must build from reviewed dependencies without
downloading unreviewed source. \textbf{Decision:}
\passthrough{\lstinline!find\_package(CONFIG)!} against approved,
versioned packages for release; allow pinned source dependencies during
development by exception. \textbf{Alternatives:} always
\passthrough{\lstinline!FetchContent!}, recursive submodules,
ExternalProject superbuild. \textbf{Verification:} disconnected CI
release job and inventory review.

\section{ADR-003: Minimum CMake
version}\label{adr-003-minimum-cmake-version}

\textbf{Context:} file sets and workflow presets are used.
\textbf{Decision:} require CMake 3.25+, not the oldest CMake supporting
only some features. \textbf{Trade-off:} fewer legacy environments,
simpler scripts. \textbf{Review trigger:} support matrix or feature
requirements change.

\section{ADR-004: Testing boundary}\label{adr-004-testing-boundary}

\textbf{Context:} in-tree tests pass, but SDK consumers report missing
symbols and headers. \textbf{Decision:} release gate runs a separate
project against only staged installation. \textbf{Consequences:} extra
setup time but stronger package-contract verification.

\chapter{Threat and trust model}\label{threat-and-trust-model}

\begin{longtable}[]{@{}
  >{\raggedright\arraybackslash}p{(\columnwidth - 6\tabcolsep) * \real{0.2500}}
  >{\raggedright\arraybackslash}p{(\columnwidth - 6\tabcolsep) * \real{0.2500}}
  >{\raggedright\arraybackslash}p{(\columnwidth - 6\tabcolsep) * \real{0.2500}}
  >{\raggedright\arraybackslash}p{(\columnwidth - 6\tabcolsep) * \real{0.2500}}@{}}
\toprule\noalign{}
\begin{minipage}[b]{\linewidth}\raggedright
Phase
\end{minipage} & \begin{minipage}[b]{\linewidth}\raggedright
Trust boundary
\end{minipage} & \begin{minipage}[b]{\linewidth}\raggedright
Example threat
\end{minipage} & \begin{minipage}[b]{\linewidth}\raggedright
Mitigation and evidence
\end{minipage} \\
\midrule\noalign{}
\endhead
\bottomrule\noalign{}
\endlastfoot
Dependency resolution & external registry -\textgreater{} CI runner &
substituted malicious package & hash/signature validation, allowlist,
access logs \\
Configure & untrusted CMake scripts -\textgreater{} shell/process &
arbitrary command execution & controlled runner, no production secrets,
reviewed changes \\
Compile/link & compiler/toolchain -\textgreater{} artifact & compromised
toolchain & pinned toolchain images, provenance \\
Test & fixture -\textgreater{} network/credentials & test exfiltrates
secrets & least-privilege test identities, restricted network \\
Package & install tree -\textgreater{} distributable & extra binaries or
licenses omitted & manifest validation, SBOM reconciliation \\
Publish & release workflow -\textgreater{} registry & unsigned or
unverified artifact & protected environment, signing, verification \\
\end{longtable}

\textbf{Responsibility split:} CMake builds the target and installation
graph; controls such as artifact attestations, vulnerability disclosure,
binary signing, provenance verification and policy enforcement are
supplied by adjacent tools and organizational systems.

\chapter{Architecture review
checklist}\label{architecture-review-checklist}

\begin{itemize}
\tightlist
\item[$\square$]
  Every target has explicit purpose, owner, API and test coverage.
\item[$\square$]
  Include paths and link requirements propagate only where they are
  needed.
\item[$\square$]
  A child project cannot unexpectedly replace the parent
  compiler/toolchain.
\item[$\square$]
  Public interfaces and installation contracts are separately verified.
\item[$\square$]
  Dependency source, integrity and licenses are recorded.
\item[$\square$]
  Configure scripts execute under least privilege with no release
  credentials.
\item[$\square$]
  Presets reflect real supported configurations and do not contain
  secrets.
\item[$\square$]
  An ADR records normal build versus superbuild selection.
\end{itemize}

\chapter{References}\label{references}

Textbook: Chapters 16, 18--19, 35, 38--43. CMake package guidance:
\url{https://cmake.org/cmake/help/latest/guide/importing-exporting/} .
This architecture review framework is an engineering extension to the
book.

\end{document}
